\section{Height, floor exclusion, and quantitative distortion}
\label{es:section}

Throughout this section the prepared strict exterior and
$0<\epsilon<1$ are fixed.  We use the filled manifolds, coefficients, and
equations of Section~\ref{sy:section}, with dimension three and $k=2$.
Thus
\[
 \ell=e^{-N\epsilon},\qquad \eta=\ell^{3/2},\qquad
 Z=t+\tfrac14\log D,\qquad
 u=e^t l\sqrt D=l e^{2Z-t}.
\]
All assertions concern smooth solutions on the truncation at $r=R$, with
$f=Z=0$ on that sphere.  They apply to the homotopy
\begin{equation}\label{es:homotopy}
 \begin{gathered}
 (a,b)\in([0,1]\times\{1\})\cup(\{0\}\times[0,1]),\\
 K_a=aK,\qquad F=h=\eta f,\qquad
 \operatorname{div}V=b u\Xi,
 \end{gathered}
\end{equation}
where all quantities involving the tensor are recomputed using $K_a$.
In particular,
\begin{equation}\label{es:source}
 \Xi=\delta_0(\mathcal T+G)+\rho
       -m_0 C_N(\rho_0+v\rho_1),\qquad
 G=\eta a_w|df|_g
   =\frac{\eta}{l}(\sqrt D-D^{-1/2}).
\end{equation}
Here $\rho_0=r^{-3-\beta}$, $\rho_1=r^{-2b_1}$,
$3/4<b_1<1$, $\rho=c_*\rho_0$, and
$m_0=\vartheta(N(t+\epsilon))$.

Constants denoted by $C$ can depend on the fixed prepared exterior and on
$\epsilon$, but not on $N$, $R$, or the homotopy parameters.
A symbol $P_N$ denotes a positive polynomial in $1+N$ with those allowed
dependencies; it can be enlarged at successive occurrences.  We will
keep polynomial bounds and bounds of the form $\exp(P_N)$ distinct.
The enlarged core has volume at most $C(1+N)$, uniformly bounded local
geometry, and uniformly bounded $K_a$ and its derivatives.  The tensor
vanishes outside that core and a fixed exterior annulus.  Radius weights
are uniformly comparable to one on the added necks and caps.

\subsection{The height and the outer face}

\begin{lemma}[Uniform height and outer slope]\label{es:height}
There are constants $C,r_0$, independent of $N$, such that every smooth
solution of the trace equation in \eqref{es:homotopy} satisfies
\begin{equation}\label{es:height-bound}
 |h|\le C,
 \qquad
 |h|\le C(r^{-b_1}-R^{-b_1})\quad(r_0\le r\le R),
\end{equation}
provided $R\ge2r_0$.  On the outer sphere,
\begin{equation}\label{es:face-slope}
 |df|_g\le C\eta^{-1}R^{-b_1-1}.
\end{equation}
Consequently, for all $R\ge R_{\min}(N,\epsilon)$, one has
$t>-\epsilon$ on that sphere.  If $a=0$, then $f=h=0$.
\end{lemma}

\begin{proof}
At an interior maximum of $h$, its gradient vanishes, and hence
$D=1$, $A_\chi=I$.  The trace equation there reads
\[
 h=\operatorname{tr}_gK_a+\frac l\eta\Delta_g h
       \le\operatorname{tr}_gK_a.
\]
At an interior minimum the reverse inequality holds.  The boundary
height is zero and $|\operatorname{tr}_gK_a|\le C$ uniformly on the
filled manifolds.  This proves the first assertion.  It also proves
$h=0$ when $K_a=0$.

Choose $r_0$ outside the support of $K$ and large enough for the
following radial calculation.  At a contact point with
$H_R=C_0(r^{-b_1}-R^{-b_1})$, the slope axis of $h$ is the radial axis.
The transverse eigenvalues of $A_\chi$ are one and its axial eigenvalue
is $0<\chi\le1$.  The asymptotic bounds for the prepared metric imply,
uniformly for this entire range of $\chi$,
\begin{equation}\label{es:radial-trace}
 \operatorname{tr}_{A_\chi}\nabla^2 H_R
 =C_0 b_1 r^{-b_1-2}
       \bigl(-2+(1+b_1)\chi+o(1)\bigr)<0.
\end{equation}
Indeed, in the Euclidean metric the radial and transverse Hessian
eigenvalues are respectively
$C_0b_1(1+b_1)r^{-b_1-2}$ and
$-C_0b_1r^{-b_1-2}$; the metric and connection errors are
$o(r^{-b_1-2})$.  Since $1-b_1>0$, one choice of $r_0$ makes the
inequality strict for every slope and every $N$.
Choose $C_0$ so that $H_R\ge C$ at $r=r_0$ for every $R\ge2r_0$.
If $h-H_R$ had a positive interior maximum, the matching gradients
would determine the same $t$, $l$, $D$, and $A_\chi$ in evaluating the
trace at that point, while $\nabla^2h\le\nabla^2H_R$.
The trace equation, $K=0$, and \eqref{es:radial-trace} would give
$0<h\le(l/(\eta\sqrt D))
\operatorname{tr}_{A_\chi}\nabla^2H_R<0$.
Applying the same argument to $-h$ proves the end bound.
The two barriers vanish at $r=R$; their one-sided derivatives there,
together with the zero tangential derivative of $h$, imply
\eqref{es:face-slope}.

For a fixed boundary slope the function
\[
 t\longmapsto t+\tfrac14\log(1+l(t)^2|df|_g^2)
\]
has derivative $1+pv/2\ge1$.  At $t=-\epsilon$ its value is at most
\[
 -\epsilon+\tfrac14\log
       (1+C^2\ell^2\eta^{-2}R^{-2b_1-2})
 =-\epsilon+\tfrac14\log
       (1+C^2\ell^{-1}R^{-2b_1-2}).
\]
This is negative for sufficiently large $R$ at fixed $N$; for example
it suffices to require
$C^2\ell^{-1}R^{-2b_1-2}<e^{4\epsilon}-1$.
Since its value at the actual boundary $t$ is $Z=0$, strict
monotonicity gives $t>-\epsilon$.
\end{proof}

\subsection{Exclusion of the lower boundary of the admissible set}

\begin{lemma}[Floor exclusion]\label{es:floor}
One can choose a fixed $\delta_0>0$ and a polynomial $C_N$ such that
no smooth solution of \eqref{es:homotopy} with $t\ge-\epsilon$
touches $t=-\epsilon$, provided $R\ge R_{\min}(N,\epsilon)$.
The choices are uniform in the homotopy parameters.
\end{lemma}

\begin{proof}
The outer boundary is excluded by Lemma~\ref{es:height}.
Consider first $b=1$ and an interior minimum with $t=-\epsilon$.
Write $e$ for the slope axis, arbitrary at zero slope.
The stationary-point calculation of
Lemma~\ref{sy:stationary-minimum} gives
\begin{align}
 \tfrac12e^{2t}R_{\widehat g}
 &=\tfrac12R_g-v\operatorname{Ric}_g(e,e)
       +\mathcal S(H^f)
       -vp\operatorname{tr}_{e^\perp}\nabla^2t
       -2\operatorname{tr}_{A_d}\nabla^2t
 \notag\\
 &\le\tfrac12R_g-v\operatorname{Ric}_g(e,e)
       +\mathcal S(H^f),                                      \label{es:min-gauss}
\end{align}
where
$\mathcal S(B)=\tfrac12((\operatorname{tr}_{A_d}B)^2
-|B|^2_{A_d\otimes A_d})$.
The signs in this inequality use $p,v\ge0$, $A_d>0$, and
$\nabla^2t\ge0$.
The same lemma supplies a constant $c_1>0$, independent of $N$, and
the algebraic estimate
\begin{equation}\label{es:min-subtraction}
 \mathcal T-\mathcal S(H^f)
       \ge c_1\mathcal T-P_N(vF^2+|K_a|^2).
\end{equation}
Both formulas concern the actual tensor $K_a$; no energy condition on
the filling or on the scaled tensor is used.

By the scalar identity \eqref{sy:scalar-equation} and $F=\eta f$,
$w(F)=G$, so comparison with \eqref{es:min-gauss} yields
\begin{align*}
 u^{-1}\operatorname{div}V
 &\ge \mathcal T-\mathcal S(H^f)+G
       +(\mu_a-R_g/2)+J_a(w)-F\operatorname{tr}_gK_a
       +v\operatorname{Ric}_g(e,e)\\
 &\ge c_1\mathcal T+G-B_N(\rho_0+v\rho_1)
\end{align*}
with $B_N\le C(1+N)$, using the coefficient $6N+25$ in
Lemma~\ref{sy:stationary-minimum}.  To justify the last bound throughout the
noncompact range, observe that
$\mu_a-R_g/2=((\operatorname{tr}K_a)^2-|K_a|^2)/2$, and that this
term, $J_a$, $K_a$, and $F\operatorname{tr}K_a$ are supported in the
enlarged core, where they are bounded and $\rho_0$ is bounded below.
On the end, $|F|^2=|h|^2\le C\rho_1$ by
Lemma~\ref{es:height}, while
$|\operatorname{Ric}_g|\le C r^{-3}\le C\rho_1$.
The same estimates on a fixed connecting annulus are absorbed into
the constant.  This proves the displayed bound with constants
independent of $R$ and $a$.

At the assumed minimum $m_0=1$.  Equating the last lower bound with
\eqref{es:source} gives
\begin{equation}\label{es:floor-contradiction}
 0\ge(c_1-\delta_0)\mathcal T+(1-\delta_0)G
       +(C_N-B_N)(\rho_0+v\rho_1)-c_*\rho_0.
\end{equation}
Choose
$0<\delta_0<\frac12\min(c_1,1)$, and then choose the polynomial
$C_N\ge B_N+c_*+1$.
The right side of \eqref{es:floor-contradiction} is positive because
$\rho_0>0$, a contradiction.

In the remaining homotopy segment $a=0$,
Lemma~\ref{es:height} gives $f=0$.  Consequently
$t=Z$, $v=0$, $A_\chi=I$,
$u=e^t l(t)$, and $\mathcal T=(2p+1)|dt|^2$.
At an interior minimum $t=-\epsilon$, the equation is
\[
 2u\Delta_g t=b u(\rho-C_N\rho_0).
\]
Its left side is nonnegative, whereas its right side is negative if
$b>0$.  If $b=0$, set
$\Phi_0(t)=\int_0^t e^s l(s)\,ds$.
Then $\Delta_g\Phi_0(t)=0$ with zero boundary data, hence
$\Phi_0(t)=0$ and $t=0$, since $\Phi_0$ is strictly increasing.
This proves the assertion on the whole homotopy.
\end{proof}

\subsection{A weighted integral above high levels}

For the rest of the section assume $t\ge-\epsilon$.  The definition
of $l$, including its constant value once $t\ge1/N$, gives
\begin{equation}\label{es:warp-bounds}
 \ell\le l\le e,\qquad t\le Z.
\end{equation}
Let $\mathcal C_N$ be the enlarged core together with a fixed annulus
containing the support of $K$.  Enlarge it by a fixed amount when
necessary.  It has volume $O(1+N)$, $\rho\ge c>0$ there, and it has
a cover by $O(1+N)$ patches of fixed size with uniform local geometry.
These properties also hold for a fixed enlargement of every patch.
Write
\[
 d\nu=\sqrt D\,dV_g,\qquad A=A_\chi.
\]
This measure is allowed to depend on the particular solution.

\begin{lemma}[High-level integral]\label{es:high-integral}
On the first segment $b=1$ of \eqref{es:homotopy}, there is a
polynomial $Z_0=P_N\ge2$ such that
\begin{equation}\label{es:high-sign}
 \Xi\ge\delta_0\mathcal T+\rho+\tfrac12\delta_0G
       \qquad\hbox{where }Z>Z_0.
\end{equation}
Every fixed enlargement of $\mathcal C_N$ satisfies
\begin{equation}\label{es:moment}
 \int_{\mathcal C_N} e^Z\,d\nu\le\exp(P_N).
\end{equation}
The constants are independent of $R$ and $a$.
\end{lemma}

\begin{proof}
Only the support of the penalty needs consideration for
\eqref{es:high-sign}.  There
$-\epsilon\le t<-\epsilon+1/N\le1$, and therefore, when $Z\ge2$,
\[
 G=\frac\eta l\left(e^{2(Z-t)}-e^{-2(Z-t)}\right)
       \ge c\eta e^{2Z}.
\]
The radius weights are bounded above, so the negative penalty has
magnitude at most $C C_N$.  Thus one may take
\[
 Z_0\ge2+\tfrac34N\epsilon+
       \tfrac12\log(1+C C_N/\delta_0).
\]
Increasing this expression to a polynomial proves
\eqref{es:high-sign}.

Choose a nondecreasing Lipschitz function $q$ which is zero on
$(-\infty,Z_0]$, one on $[Z_0+1,\infty)$, and has $0\le q'\le2$.
It vanishes on the outer boundary.  Testing the divergence equation
with $q(Z)$ gives
\begin{align}
 \int u\Xi q(Z)\,dV_g
 &=-\int u q'(Z)
       \bigl(2|dZ|_A^2+A(K_a(w,\cdot),dZ)\bigr)\,dV_g
 \notag\\
 &\le C\int_{\{Z_0<Z<Z_0+1\}}u|K_a|^2\,dV_g.       \label{es:cutoff-test}
\end{align}
The last inequality is the scalar estimate
$-2x^2+yx\le y^2/8$, and $A\le I$.
The last integrand is supported in $\mathcal C_N$, where
$u=l e^{2Z-t}\le e^{1+\epsilon+2(Z_0+1)}$ on the strip.
Its integral is consequently at most $\exp(P_N)$.
All terms on the left are nonnegative by \eqref{es:high-sign}, and
$q=1$ above $Z_0+1$, so
\begin{equation}\label{es:weighted-positive}
 \int_{\{Z>Z_0+1\}}u(\delta_0\mathcal T+\rho+
                   \tfrac12\delta_0G)\,dV_g
       \le\exp(P_N).
\end{equation}
The same argument applies to any fixed enlargement of the core,
since $\rho$ has a positive lower bound there.

Here is a pointwise conversion of this estimate to the claimed
moment.  Since $l\ge\ell\ge\eta$ and $D\ge1$,
\begin{align*}
 u(1+G)
 &=\sqrt D e^t\bigl[l+\eta(\sqrt D-D^{-1/2})\bigr]\\
 &\ge\eta e^t D
 \ge\eta e^Z\sqrt D.
\end{align*}
The final inequality is $\sqrt D=e^{2(Z-t)}\ge e^{Z-t}$.
Thus \eqref{es:weighted-positive}, $\rho\ge c$ on the core, and
$\eta^{-1}=e^{3N\epsilon/2}$ bound the integral in
\eqref{es:moment} above $Z_0+1$.
On the complementary set,
$\sqrt D=e^{2(Z-t)}\le e^{2(Z_0+1+\epsilon)}$ and
$e^Z\le e^{Z_0+1}$; multiplication by the core volume
$C(1+N)$ gives the same bound there.
\end{proof}

\subsection{Sobolev and energy estimates on graph patches}

\begin{lemma}[Sobolev inequality on the unit-warp graph]
\label{es:sobolev}
For each of the core patches just described, and every smooth
compactly supported function $\varphi$ in the patch,
\begin{equation}\label{es:sobolev-form}
 \left(\int|\varphi|^6\,d\nu\right)^{1/3}
 \le S_N\int\bigl(|d\varphi|_A^2
                     +(1+\mathcal T)\varphi^2\bigr)\,d\nu,
 \qquad \log S_N\le P_N.
\end{equation}
No bound on $Z$, $df$, or derivatives of $t$ is required beyond
$t\ge-\epsilon$ and the trace and coercivity identities.
\end{lemma}

\begin{proof}
Use the graph of $f$ in the product metric $g+dz^2$ on a larger base
patch times the entire real line.  Put
\[
 E_f=1+|df|_g^2,\qquad
 A_1=I-\frac{df\otimes df}{E_f}.
\]
Its measure is $d\nu_1=\sqrt{E_f}\,dV_g$ and its inverse metric on
base covectors is $A_1$.  Equations~\eqref{es:warp-bounds} give
\[
 e^{-2}D\le E_f\le\ell^{-2}D.
\]
The axial eigenvalues of $A_1,A_d,A$ are respectively
$E_f^{-1},d,\chi$, and
$2d/(2+N)\le\chi\le d$.  Hence the measures and inverse metrics
are pairwise comparable with factors bounded by
$C(2+N)\ell^{-2}\le\exp(P_N)$.

The scalar mean curvature of this graph in $g+dz^2$ is
\[
 H_1=E_f^{-1/2}\operatorname{tr}_{A_1}\nabla^2f
    =\frac{\sqrt D}{l\sqrt{E_f}}
           \operatorname{tr}_{A_1}H^f.
\]
Coercivity, Proposition~\ref{sy:coercivity}, controls the transverse block
of $H^f+K_a$ and the axial block multiplied by $\sqrt d$ by
$P_N\sqrt{\mathcal T}$.  Since
$E_f^{-1}\le e^2d\le e^2\sqrt d$ and $|K_a|\le C$, it follows that
\[
 |\operatorname{tr}_{A_1}H^f|
       \le P_N(1+\sqrt{\mathcal T}),\qquad
 |H_1|\le\exp(P_N)(1+\sqrt{\mathcal T}).
\]
This estimate uses weighted axial contraction, rather than an
unweighted bound for the full Hessian.

Each larger base patch can be smoothly isometrically embedded in a
fixed Euclidean space.  Only finitely many compact geometric models
are needed: the fixed exterior and transition patches, the fixed
caps, and translates of product-neck patches.  One can obtain such
embeddings by extending each model to a compact smooth Riemannian
manifold and applying the smooth isometric embedding theorem
\cite[Theorem~2, p.~59]{Nash1956}.
The Euclidean dimensions and the second fundamental forms on the
smaller patches therefore have fixed bounds.  Taking the product
of these embeddings with the identity of the real line gives an
isometric embedding of the ambient product.  The Euclidean mean
curvature vector of the graph has norm at most $|H_1|+C$, because
the trace of the ambient second fundamental form over its three
unit tangent vectors is bounded by a fixed constant.

The Euclidean submanifold Sobolev inequality of Michael and Simon
\cite{MichaelSimon1973}, in the form proved in
\cite[Chapter~4, Section~6, Theorem~6.7]{Simon2018GMT},
applies to this smooth three-dimensional immersed graph and compactly
supported test functions: if $\zeta\ge0$, then
\[
 \left(\int\zeta^{3/2}\,d\nu_1\right)^{2/3}
 \le C\int\bigl(|d\zeta|_{A_1}+|\mathbf H|\zeta\bigr)\,d\nu_1.
\]
There is no boundary term because the support is in the interior of
the larger patch.  Apply it to $\zeta=|\varphi|^4$ and use
Cauchy--Schwarz on
$4|\varphi|^3|d\varphi|$ and
$|\mathbf H||\varphi|^4$ to obtain
\[
 \left(\int|\varphi|^6\,d\nu_1\right)^{1/3}
 \le C\int\bigl(|d\varphi|_{A_1}^2+
                          |\mathbf H|^2\varphi^2\bigr)\,d\nu_1.
\]
The zero function is harmless; otherwise this follows by dividing
the intermediate inequality by
$(\int|\varphi|^6\,d\nu_1)^{1/2}$ and squaring.
The already established metric, measure, and mean curvature
comparisons yield \eqref{es:sobolev-form}.
\end{proof}

\begin{lemma}[Normalized energy estimate]\label{es:energy}
There are polynomials $E_N$ and $s_0(N)\ge1$ such that on the first
homotopy segment, for $s\ge s_0(N)$ and any compactly supported base
cutoff $\psi$,
\begin{equation}\label{es:energy-form}
 \int\psi^2e^{2sZ}\bigl(\mathcal T+s|dZ|_A^2\bigr)\,d\nu
 \le E_N(1+s)\int e^{2sZ}
                     (\psi^2+|d\psi|_g^2)\,d\nu.
\end{equation}
All constants are independent of $R$, $a$, and the solution.
\end{lemma}

\begin{proof}
Put $B=e^t l$, $\kappa=K_a(w,\cdot)$ and
$\gamma=d\log B=(1+p)dt$.  Coercivity and $A\le A_d$ show that
\begin{equation}\label{es:logweight}
 |\gamma|_A\le B_N\sqrt{\mathcal T},\qquad
 |\kappa|_A\le C,
\end{equation}
with $B_N$ polynomial.  The boundedness of the weights in the
penalty gives, globally,
$\Xi\ge\delta_0\mathcal T-L_N$ for a polynomial $L_N$.
Test $\operatorname{div}V=u\Xi$ with the nonnegative function
$\psi^2e^{2sZ}/B$.  As $u/B=\sqrt D$, exact differentiation gives
\begin{align}
 &\int\psi^2e^{2sZ}\Xi\,d\nu
       +4s\int\psi^2e^{2sZ}|dZ|_A^2\,d\nu\notag\\
 &=\int e^{2sZ}\bigl[
       \psi^2 A(\gamma,2dZ+\kappa)
       -2\psi A(d\psi,2dZ+\kappa)
       -2s\psi^2A(\kappa,dZ)\bigr]\,d\nu.             \label{es:normalized-test}
\end{align}
For clarity, set $X=|dZ|_A$ and $Y=|d\psi|_A\le|d\psi|_g$.
Young's inequality and \eqref{es:logweight} give the following
pointwise bounds for the terms on the right:
\begin{align*}
 2B_N\psi^2\sqrt{\mathcal T}X
     &\le\tfrac14\delta_0\psi^2\mathcal T
                         +C\delta_0^{-1}B_N^2\psi^2X^2,\\
 C B_N\psi^2\sqrt{\mathcal T}
     &\le\tfrac14\delta_0\psi^2\mathcal T
                         +C\delta_0^{-1}B_N^2\psi^2,\\
 4|\psi|YX&\le s\psi^2X^2+4s^{-1}Y^2,\\
 2C|\psi|Y&\le C(\psi^2+Y^2),\\
 2Cs\psi^2X&\le s\psi^2X^2+C s\psi^2.
\end{align*}
Choose the polynomial threshold
$s_0\ge1+C\delta_0^{-1}B_N^2$, with its constant large enough
that the first $X^2$ error is at most $s\psi^2X^2$.
Substituting these five estimates into
\eqref{es:normalized-test} leaves at least
$\delta_0\mathcal T/2+sX^2$ on its left and at most
$P_N(1+s)(\psi^2+Y^2)$ on its right, all multiplied by
$e^{2sZ}$.  Absorbing the fixed factor $\delta_0/2$ proves
\eqref{es:energy-form}.
\end{proof}

\subsection{Iteration with a polynomial logarithmic bound}

\begin{proposition}[Quantitative distortion]\label{es:distortion}
For the choices in Lemma~\ref{es:floor}, every smooth solution of
\eqref{es:homotopy} with $t\ge-\epsilon$ and
$R\ge R_{\min}(N,\epsilon)$ satisfies
\begin{equation}\label{es:distortion-bound}
 -\epsilon<t\le Z\le P_N,\qquad
 \ell\le l\le e,\qquad
 |f|+|df|_g\le\exp(P_N).
\end{equation}
The polynomial is independent of the outer radius and the homotopy
parameters.  In particular, exhausting the exterior at fixed $N$
preserves these bounds.
\end{proposition}

\begin{proof}
First take $b=1$ and write $Y=e^Z$.  Choose nested core patches
$U_r\Subset U_R$ with $1/2\le r<R\le1$, in uniform local
coordinates, and cutoffs with $|d\psi|_g\le C/(R-r)$.
Combining Lemmas~\ref{es:sobolev} and~\ref{es:energy}, applied to
$\varphi=\psi e^{sZ}$, gives
\begin{equation}\label{es:one-step}
 \|Y\|_{L^{6s}(U_r,d\nu)}
 \le\bigl[C S_NE_N(1+s)^2(1+(R-r)^{-2})\bigr]^{1/(2s)}
        \|Y\|_{L^{2s}(U_R,d\nu)}.
\end{equation}
Indeed, differentiation of $\psi e^{sZ}$ contributes at most
$2e^{2sZ}|d\psi|_A^2+2s^2\psi^2e^{2sZ}|dZ|_A^2$; the latter
integral is bounded by $2s$ times \eqref{es:energy-form}.
This accounts for the power $(1+s)^2$ in
\eqref{es:one-step}.

Set $s_i=3^i s_0$, take successive spatial gaps proportional to
$b2^{-i-1}$, where $0<b\le1$ is the total gap, and iterate
\eqref{es:one-step}.  The product of constants has logarithm at most
\[
 \sum_{i=0}^{\infty}\frac{
  \log(C S_NE_N)+2\log(1+s_0)+C i+2\log(1/b)}{2s_0 3^i}.
\]
The identities
\[
 \sum_{i\ge0}3^{-i}=\tfrac32,
 \qquad \sum_{i\ge0}i3^{-i}=\tfrac34
\]
therefore give
\begin{equation}\label{es:local-moser}
 \log\sup_{U_r}Y
 \le\frac1{2s_0}\log\int_{U_R}Y^{2s_0}\,d\nu
       +\frac{D_N}{s_0}
       +\frac{C}{s_0}\log\frac1{R-r},
\end{equation}
where
$D_N\le C(1+\log S_N+\log E_N+\log(1+s_0))\le P_N$.
The limiting norm is the supremum because each individual solution
and its positive smooth measure are smooth on the closure of the
larger patch.  No uniform bound for that individual supremum has
been used to obtain \eqref{es:local-moser}.

We next lower the initial integrability exponent explicitly.
By Lemma~\ref{es:high-integral},
$\int_{U_R}Y\,d\nu\le e^{L_N}$ with $L_N$ polynomial.
Put $M(r)=\max(1,\sup_{U_r}Y)$ and
$\theta=1-1/(2s_0)$.  The inequality
$Y^{2s_0}\le M(R)^{2s_0-1}Y$ and
\eqref{es:local-moser} yield
\begin{equation}\label{es:interpolation}
 \log M(r)\le\theta\log M(R)
       +\frac{L_N}{2s_0}+\frac{D_N}{s_0}
       +\frac C{s_0}\log\frac1{R-r}.
\end{equation}
Choose increasing radii $r_j$ whose successive gaps are
$c2^{-j-1}$ and whose limit lies strictly inside the larger patch.
Iterating \eqref{es:interpolation} $m$ times gives
\[
 \log M(r_0)
 \le\theta^m\log M(r_m)
       +\sum_{j=0}^{m-1}\theta^j
           \left(\frac{L_N/2+D_N+C}{s_0}
                         +\frac{Cj}{s_0}\right).
\]
The first term tends to zero: the supremum on that fixed larger
patch is finite for the particular smooth solution.
Moreover,
\[
 \sum_{j\ge0}\theta^j=2s_0,
 \qquad \sum_{j\ge0}j\theta^j
             =\frac{\theta}{(1-\theta)^2}\le4s_0^2.
\]
Consequently
\begin{equation}\label{es:log-bound}
 \log M(r_0)\le C(L_N+D_N+1+s_0)\le P_N.
\end{equation}
This computation is the required control of the parameter dependence:
both iteration steps enlarge only a polynomial in the logarithm of
the supremum.  Uniform patches cover the entire enlarged core and
the support of $K$, so \eqref{es:log-bound} bounds $Z$ there with
one polynomial independent of the number of patches.

Outside this covered region $K=0$.  If $Z$ attained a larger
maximum, exceeding also $Z_0$, the maximum would be interior since
$Z=0$ on the outer sphere.  At such a point
$V=2uA\nabla Z$ and
\[
 \operatorname{div}V=2u\operatorname{tr}_A\nabla^2Z\le0,
 \qquad u\Xi\ge u\rho>0
\]
by \eqref{es:high-sign}, a contradiction.  This proves the global
upper bound on the first segment.

On the second segment $a=0$, the trace equation again gives $f=0$
and $t=Z$.  If $b>0$, a positive interior maximum with
$Z>\max(0,-\epsilon+1/N)$ has $m_0=0$, $dt=0$, and hence
$\operatorname{div}V\le0<b u\rho$.
Together with zero boundary values this is already a uniform upper
bound.  If $b=0$, the transformed harmonic equation used in
Lemma~\ref{es:floor} gives $Z=0$.

Finally Lemma~\ref{es:floor} makes the lower bound strict,
\eqref{es:warp-bounds} gives the warp bounds, and
$D=e^{4(Z-t)}\le\exp(P_N)$ gives
$|df|_g=\sqrt{D-1}/l\le\exp(P_N)$.
Lemma~\ref{es:height} and $\eta^{-1}=e^{3N\epsilon/2}$ give
$|f|\le\exp(P_N)$.  This proves
\eqref{es:distortion-bound}.
\end{proof}

The estimates above use only the extended metric and tensor in the
filled region.  The electromagnetic flux forms are needed solely on
the original exterior, and none has been extended through a cap.
The quantitative conclusions concern height, slope, and distortion;
subsequent elliptic estimates at a fixed $N$ may have arbitrary
dependence on that fixed parameter.
